\bookchapter{第五章：讓手機更懂你——基礎設定}

\chapterintro{花幾分鐘設定，體驗提升巨大。這些設定就像給手機“調音”，讓它更符合你的使用習慣。}

\booksection{一、聲音與震動}

設個喜歡的鈴聲和通知聲。靈活使用靜音模式（開會）、勿擾模式（睡覺），世界瞬間清淨。

\begin{itemize}
  \item \textbf{鈴聲}：選一個你能聽到的、不尷尬的。
  \item \textbf{震動模式}：靜音時靠震動提醒。
  \item \textbf{勿擾模式}：可以設定例外，比如允許通訊錄來電。
  \item \textbf{音量分離}：媒體音量、鈴聲音量、鬧鐘音量可以分開調。
\end{itemize}

\booksection{二、顯示與亮度}

\begin{itemize}
  \item \textbf{自動亮度}：讓手機根據環境光自動調節，省心。
  \item \textbf{護眼/夜覽模式}：過濾藍光，晚上看手機舒服點。
  \item \textbf{深色模式}：省電（對OLED屏特別有效），且夜間不刺眼。很多App也支援深色模式。
  \item \textbf{螢幕重新整理率}：部分安卓機可調，越高（如120Hz）滑動越絲滑。
  \item \textbf{自動鎖屏}：設定15秒或30秒，安全又省電。
  \item \textbf{字型大小}：覺得字小？調大一點。
\end{itemize}

\booksection{三、桌布、主題與桌面}

換張好看的桌布（靜態或動態），心情都變了。安卓品牌手機常有主題商店，能整套換圖示和風格。

\textbf{桌面小元件}：天氣、日程、音樂控制放桌面，一步直達。iOS和安卓都支援，長按桌面空白處新增。

\booksection{四、效率設定：讓手機成為你的“外掛大腦”}

這部分是“一旦設定就回不去”的功能，每個都能讓你的日常效率顯著提升。

\booksubsection{專注 / 免打擾模式：奪回被通知奪走的注意力}

\begin{itemize}
  \item \textbf{核心思想}：不是"少看手機"，而是"在該專注時專注，該聯絡時再聯絡"。
  \item \textbf{使用建議}：\newline · 工作模式：僅允許"星標"聯絡人和必要 App 透過；\newline · 睡眠模式：自動開啟，僅允許鬧鐘；\newline · 健身 / 閱讀 / 會議：根據當前活動定製。
  \item \textbf{進階玩法}：用自動化 + 位置 / 時間觸發，讓模式自動切換。例如：iOS 快捷指令可設定"到達公司即自動進入工作模式"。
\end{itemize}

\booksubsection{通知管理：讓紅點不再“奪命連環 call”}

\begin{itemize}
  \item \textbf{關閉營銷通知}：90\% 的推送其實你並不需要。
  \item \textbf{按重要性排序}：把人際通訊類（電話、簡訊、即時通訊）置頂；新聞、購物、社交平臺類可設為"靜默推送"。
  \item \textbf{分組通知}：把同類通知摺疊在一起，避免一次解鎖被幾十條訊息淹沒。
\end{itemize}

\booksubsection{快捷手勢：讓操作快人一步}

\begin{itemize}
  \item \textbf{截圖}：電源 + 音量下（iOS 全面屏）/電源 + 音量下（安卓通用）。
  \item \textbf{截長圖 / 滾動截圖}：普通截圖後點選"長截圖"選項。
  \item \textbf{一鍵掃碼}：控制中心 / 快捷設定裡新增"掃一掃"圖示。
  \item \textbf{輕點背面（iPhone）}：在 [設定] > [輔助功能] > [觸控] > [輕點背面] 中繫結"截圖"或"鎖屏"等動作。
  \item \textbf{語音喚醒}：Siri / 小愛同學 / 小藝 / 小布，喊一聲就能查天氣、設鬧鐘。
\end{itemize}

\booksubsection{輸入法：最被低估的效率神器}

\begin{itemize}
  \item \textbf{選對輸入法}：\newline · 搜狗 / 百度：詞庫大，流行語更新快；\newline · 訊飛：語音輸入識別率最高，開會記筆記神器；\newline · Gboard（谷歌）：簡潔、跨平臺同步強、支援滑行輸入；\newline · iOS 自帶：隱私最好，與系統深度整合。
  \item \textbf{必備開關}：智慧糾錯、聯想詞、雲詞庫同步、上滑輸入數字符號。
  \item \textbf{隱私提示}：慎用"整句聯想"或"個性化推薦"功能，可能上傳輸入內容。
\end{itemize}

\booksubsection{桌面與多工：讓工作流順暢起來}

\begin{itemize}
  \item \textbf{App 分屏 / 畫中畫}：看影片時回訊息，上網時記筆記。
  \item \textbf{小窗 / 懸浮窗}：微信 / 抖音支援"小窗模式"，回覆完即收起。
  \item \textbf{常用 App 置頂}：Dock 欄（螢幕底部）只放最高頻的 4 個，減少查詢時間。
\end{itemize}

\begin{tcbinfo}{"上工第一天"檢查清單}
拿到新手機或大版本更新後，建議花 10 分鐘過一遍這些設定：\newline (1) 鎖屏密碼 + 生物識別 (2) 通知分組 + 關閉營銷推送 (3) 字型大小 / 重新整理率 / 深色模式\newline (4) 預設輸入法 / 瀏覽器 / 郵件 (5) 緊急聯絡人 / 醫療卡 (6) 備份設定

\end{tcbinfo}

\booksection{五、輔助功能與無障礙：科技向善的力量}

別被“輔助功能”這個名字嚇到，這些功能人人能用。它們最初是為殘障人士設計的，但對普通人也很有用——臨時手受傷抱娃、地鐵上耳機壞了聽不清、戴老花鏡看小字……都能救場。說到底，無障礙設計體現了\textbf{科技向善}的理念：讓每個人都能平等地享受數字生活。

\booksubsection{視障使用者的“眼睛”——螢幕閱讀器全攻略}

螢幕閱讀器是視障使用者使用手機的核心工具。它會把螢幕上\textbf{所有可視內容朗讀出來}，使用者透過手勢和語音操作手機。下面分別介紹 iOS “旁白”和安卓 TalkBack 的使用方法。

\begin{tcbinfo}{iOS 旁白（VoiceOver）}
\textbf{開啟路徑}：[設定] > [輔助功能] > [旁白]，開啟開關。也可三擊側邊鍵（或主螢幕鍵）快速開關。

\textbf{基本手勢}：

\begin{enumerate}
  \item 單指滑動：逐元素瀏覽螢幕內容；
  \item 雙指輕掃：上下/左右整頁滾動；
  \item 單指雙擊：啟用當前焦點元素（相當於"點選"）；
  \item 三指輕掃：切換頁面；
  \item 單指三擊：放大或選擇（與普通三擊不同）。
\end{enumerate}

\textbf{個性化設定}：[旁白] > [語音]，可調節\textbf{朗讀語速}（預設約 50\%，可調到 70-80\% 聽感更流暢）、\textbf{音調}、切換不同語音包。iOS 自帶“增強語音”和“語音控制器”，前者能減少呼吸聲和停頓，聽感更自然。

\end{tcbinfo}

\begin{tcbinfo}{安卓 TalkBack}
\textbf{開啟路徑}：[設定] > [輔助功能] > [TalkBack]，開啟開關。或長按音量上下鍵 3 秒快速啟用（部分品牌支援）。

\textbf{基本手勢}：

\begin{enumerate}
  \item 單指左右滑動：切換焦點元素；
  \item 單指雙擊：啟用元素；
  \item 雙指上下/左右滑：滾動；
  \item 三指單擊：進入閱讀模式（連續朗讀）；
  \item 四指單擊：暫停/繼續朗讀。
\end{enumerate}

\textbf{個性化設定}：[TalkBack] > [文字轉語音設定]，可調整\textbf{語速}、\textbf{音調}、切換\textbf{語音引擎}（建議安裝 Google 文字轉語音引擎，朗讀更自然）。

\end{tcbinfo}

\booksubsection{讓小字變大、讓模糊變清晰——視覺調整}

\booksubsection{放大手勢}

\begin{itemize}
  \item \textbf{iOS}：[設定] > [輔助功能] > [縮放]，開啟後\textbf{三指輕點螢幕}即可放大。預設全屏縮放，可選視窗縮放（區域性放大窗）。\textbf{雙指開合可調節倍數（1.5x 至 15x）}，三指雙擊鎖定當前縮放比例。
  \item \textbf{安卓}：[設定] > [輔助功能] > [放大]，三指點選放大，再三指點選退出；或直接\textbf{三指點選螢幕三下}啟用臨時放大（最多 8 倍）。
\end{itemize}

\booksubsection{高對比度文字 / 粗體文字}

\begin{itemize}
  \item \textbf{iOS}：[設定] > [輔助功能] > [顯示與文字大小]，開啟"粗體文字"、"增強對比度"、"降低透明度"等開關。"智慧反轉"可在淺色背景下讓文字保持黑色，避免白底白字。
  \item \textbf{安卓}：[設定] > [顯示] > [字型大小]，可調節"小、普通、大、超大"四檔。華為/小米額外提供"粗體"和"高對比度文字"開關。
\end{itemize}

\booksubsection{色彩濾鏡（色盲友好）}

\begin{itemize}
  \item \textbf{iOS}：[設定] > [輔助功能] > [顯示與文字大小] > [色彩濾鏡]，提供\textbf{紅/綠/藍/黃四種色盲模式}，以及"灰度"和"反轉顏色"選項。建議先做一次色覺測試（系統自帶）再選對應濾鏡。
  \item \textbf{安卓}：[設定] > [輔助功能] > [顏色校正]，含\textbf{三種色盲模式}+ "反轉顏色"開關。
\end{itemize}

\booksubsection{AI 影象描述：讓 AI 當你的“眼睛”}

當螢幕閱讀器遇到圖片時，傳統做法是跳過或只讀出“圖片”兩個字。但現在 AI 可以\textbf{用一句話描述圖片內容}，對識別驗證碼、讀圖、識別選單非常有幫助。

\begin{tcbinfo}{各系統 AI 影象描述支援情況}
\begin{itemize}
  \item \textbf{iOS 15+}：在 [設定] > [輔助功能] > [旁白] > [語音描述] 中開啟"影象描述"。系統會呼叫\textbf{本地 AI 模型}自動生成圖片文字描述（需 iPhone 12 及以上）。
  \item \textbf{新版 TalkBack 應用}：Google TalkBack 在較新版本中整合\textbf{圖片問答}功能，可以問"這張圖裡有什麼"，AI 實時回答（不嚴格依賴安卓系統版本，更新 TalkBack 應用即可體驗）。
  \item \textbf{第三方 App}："Seeing AI"（iOS）、"Lookout"（安卓）、"小藝幫幫"（華為）都是優秀的 AI 視覺助手，可識別物品、讀檔案、識別人臉表情。
\end{itemize}

\end{tcbinfo}

\begin{tcbpractice}{使用場景}
試著對著一張街拍照片問“這張圖裡有幾個人？他們在做什麼？”，或者開啟外賣 App，讓 AI 讀出選單圖裡的菜品名——\textbf{這是真正的“科技讓生活更平等”}。

\end{tcbpractice}

\booksubsection{聽障使用者的“耳朵”——實時字幕全攻略}

實時字幕（Live Captions）是\textbf{系統級的語音轉文字}功能，聽障使用者在嘈雜環境、圖書館、外語場景下都特別有用。

\begin{tcbinfo}{實時字幕開啟與最佳化}
\textbf{iOS}：[設定] > [輔助功能] > [實時字幕]，開啟開關。系統會檢測裝置所有音訊（通話、影片、播客、FaceTime、語音備忘錄），自動顯示字幕在螢幕上方。

\textbf{安卓 12+}：按\textbf{音量鍵}點選彈出框中的“實時字幕”圖示；或在 [設定] > [輔助功能] > [實時字幕] 中開啟。

\textbf{準確率最佳化技巧}：

\begin{enumerate}
  \item 儘量在\textbf{安靜環境}下使用，遠離空調、機器轟鳴；
  \item 說話者\textbf{咬字清晰、語速適中}時識別率最高；
  \item 多人對話會降低準確率，可使用"只顯示一人"模式；
  \item 支援中文普通話、粵語、英語等主流語種，\textbf{小語種識別率較低}；
  \item 說話中途有背景音樂或人聲干擾，識別會斷斷續續。
\end{enumerate}

\end{tcbinfo}

\booksubsection{通話場景雙通道：語音轉文字 + 文字轉語音}

當使用者\textbf{既聽不到也說不了話}時（如突發失聲、喉部手術恢復期、嚴重口吃），手機還能當“傳話員”。

\begin{tcbinfo}{通話實時轉譯（TTY / RTT）}
\textbf{iOS}：[設定] > [輔助功能] > [RTT/TTY]，開啟後通話時螢幕上會顯示鍵盤，對方的語音會被實時打成文字，\textbf{你打的字也會被念出來}給對方聽。

\textbf{安卓}：[設定] > [輔助功能] > [實時轉錄] + [聲音通知]。“實時轉錄”將對方說話實時顯示為文字；“聲音通知”在不便語音時將你的輸入轉為語音播放。

\textbf{相容性}：

\begin{itemize}
  \item \textbf{RTT（Real-Time Text）}：雙方都需支援，文字邊打邊顯示，無延遲，主流運營商已支援；
  \item \textbf{TTY（電傳打字）}：需專用 TTY 裝置，已逐步被 RTT 取代；
  \item 微信/QQ 的\textbf{語音轉文字功能}也是替代方案：長按語音訊息 "轉文字"。
\end{itemize}

\end{tcbinfo}

\booksubsection{面向所有人的“小但實用”功能}

\begin{itemize}
  \item \textbf{輕點背面}（iOS 14+）：[設定] > [輔助功能] > [觸控] > [輕點背面]，可繫結\textbf{雙擊/三擊}觸發截圖、鎖屏、控制中心等動作。雙手拿快遞時特別好用。
  \item \textbf{輔助觸控 / 懸浮球}：iOS [輔助功能] > [觸控] > [輔助觸控]，安卓 [輔助功能] > [懸浮導航]。虛擬按鈕替代手勢或物理按鍵，按鍵不靈或單手操作時救場。
  \item \textbf{切換控制}：透過外部開關、單次點選、頭部動作控制手機，\textbf{為肢體障礙使用者設計}。iOS [輔助功能] > [切換控制]，安卓 [輔助功能] > [開關控制]。
  \item \textbf{聲音識別}（iOS）：[輔助功能] > [聲音識別]，開啟後手機會"聽"門鈴、嬰兒哭聲、火警警報、警笛等聲音，\textbf{識別到就震動+螢幕通知}，對聽障使用者和睡眠期間非常有價值。
  \item \textbf{LED 閃爍提醒}：iOS [輔助功能] > [音訊/視覺] > [LED 閃爍以示提醒]，安卓 [設定] > [通知] > [閃光燈提醒]。來電和通知時閃光燈閃爍，\textbf{嘈雜環境無聲也能感知}。
  \item \textbf{單聲道音訊}：[設定] > [輔助功能] > [音訊] > [單聲道音訊]，將左右聲道合併。單耳聽力受損的使用者不會丟失另一半聲音資訊。
\end{itemize}

\begin{tcbinfo}{蘋果的無障礙設計哲學}
蘋果在無障礙領域投入巨大，每年都會更新和新增輔助功能。蘋果官網有專門的無障礙頁面，詳細介紹各項功能的使用方法。iOS 的無障礙功能通常比安卓更完善、體驗更流暢，這也是很多殘障使用者選擇蘋果的重要原因。

\end{tcbinfo}

\begin{tcbpractice}{花30秒體驗一下}
找一段空閒時間，做以下體驗：

\begin{enumerate}
  \item 開啟\textbf{旁白 / TalkBack}，試著不靠眼睛用三指滑動瀏覽微信；
  \item 開啟\textbf{實時字幕}，對著手機說話看識別準確率；
  \item 用\textbf{三指點選}放大一段文字到 200\%，體驗視障使用者的日常；
  \item 關掉之前所有功能，思考"如果我永遠需要這些功能，我會希望它怎麼設計"。
\end{enumerate}

\end{tcbpractice}
